\begin{table}[htbp]
\centering
\caption{The phosphorus contrast in each of the five windows, with the multiplicity of testing five windows made explicit. Two-sided Mann--Whitney on six phosphorus-containing against three phosphorus-free agents.}
\label{tab:si-pfmult}
\footnotesize
\resizebox{\textwidth}{!}{%
\begin{tabular}{lrrrrrr}
\toprule
Window (\wn) & P-containing & No P & $\Delta$ & $p$ & Bonferroni & Holm \\
\midrule
250-600 skeletal/torsion & 0.810 & 0.675 & -0.135 & 0.117 & 0.583 & 0.467 \\
600-850 & 0.636 & 0.912 & +0.276 & 0.117 & 0.583 & 0.467 \\
850-1150 & 0.840 & 0.793 & -0.047 & 0.517 & 1.000 & 0.767 \\
1150-1400 & 0.814 & 0.755 & -0.059 & 0.383 & 1.000 & 0.767 \\
1400-1800 CH deform & 0.758 & 0.557 & -0.201 & 0.017 & 0.083 & 0.083 \\
\bottomrule
\end{tabular}
}

\begin{flushleft}\footnotesize No window survives correction. The smallest uncorrected $p$ is not the 600--850\,\wn{} window but 1400--1800\,\wn{}, where the ordering is reversed. With six against three molecules the smallest two-sided $p$ the design can produce is 0.024, so this test has essentially no power and we do not quote a $p$ value in support of the localisation. The argument rests on the effect size, $\Delta = +0.276$ in the 600--850\,\wn{} window against $|\Delta| \leq 0.201$ elsewhere, and on the sign reversal in all four other windows.\end{flushleft}
\end{table}
